% !TEX root = main.tex
\section{Теоретическая часть}

Для квадратных матриц $A$ и $B$ порядка $N$ элемент матрицы $C=AB$ вычисляется по формуле
\begin{equation}
c_{ij}=\sum_{k=0}^{N-1} a_{ik}b_{kj},\qquad 0\le i,j<N,
\label{eq:product}
\end{equation}
где $i$ и $j$~--- номера строки и столбца результата, а $k$~--- индекс суммирования. Для вычисления одной строки $C$ требуется одна строка $A$ и вся матрица $B$. Всего выполняется $N^3$ умножений и столько же сложений с аккумулятором, начальное значение которого равно нулю.

В OpenMP строки распределяются между потоками, которые имеют доступ к общей памяти~\cite{openmp}. В MPI каждый процесс работает со своими массивами и получает необходимые данные через операции обмена~\cite{mpi}. В гибридном варианте строки распределяются сначала между процессами MPI, затем между потоками OpenMP внутри каждого процесса.

\section{Ход работы}

\subsection{Подготовка данных}

Матрицы хранятся в одномерных массивах типа \code{int32_t}. Исходные элементы получены генератором xorshift32 с начальным состоянием 2026 и лежат в диапазоне от $-9$ до $9$. Матрица $B$ сразу записывается в транспонированном виде. Поэтому при вычислении суммы~\eqref{eq:product} по строкам $A$ и $B^\mathsf{T}$ обращения к обоим массивам идут последовательно. При $N\le4096$ модуль промежуточной суммы не превышает $81\cdot4096=331\,776$ и представим типом \code{int32_t}.

\clearpage
\subsection{Последовательный вариант и OpenMP}

Все варианты собраны из одного файла \code{src/matmul.cpp}. В листинге~\ref{lst:kernel} приведено общее ядро. В последовательной сборке внешний цикл проходит по строкам по порядку. При сборке с OpenMP используется \code{schedule(static)}, поскольку на каждую строку приходится одинаковый объём работы. Каждый поток записывает только выделенные ему строки результата.

\subsection{MPI и гибридный вариант}

В коммуникаторе \code{MPI_COMM_WORLD} ранги нумеруются от 0 до $P-1$, где $P$~--- число процессов~\cite{mpi}. На ранге 0 формируются входные матрицы и выделяется память под итоговый массив. При разбиении $N=Pq+r$, где $q$~--- целая часть от деления $N$ на $P$, а $r$~--- остаток, первые $r$ рангов получают по $q+1$ строке, остальные~--- по $q$.

Распределение данных приведено в листинге~\ref{lst:mpi}. Вызовом \code{MPI_Bcast} рассылается вся матрица $B^\mathsf{T}$. Через \code{MPI_Scatterv} каждый ранг получает свой блок строк $A$, а через \code{MPI_Gatherv} готовые строки собираются на ранге 0 в исходном порядке. В гибридной сборке внутри каждого ранга используется схема OpenMP из листинга~\ref{lst:kernel}. Для неё запрашивается уровень \code{MPI_THREAD_FUNNELED}, поскольку вызовы MPI выполняет поток, инициализировавший MPI.

Для двух рангов объём передаваемых элементов в байтах равен $4N^2+8N\lfloor N/2\rfloor$. Первое слагаемое соответствует матрице $B^\mathsf{T}$, второе~--- передаче удалённому рангу строк $A$ и возврату такого же числа строк $C$. При $N=2048$ получается 32~МиБ полезных данных без учёта служебного трафика TCP.
